% Сборка: lualatex lab4-report.tex (рисунки — node ../clip-lab/scripts/report-figures.cjs)
\documentclass[11pt,a4paper]{article}
\usepackage[a4paper,margin=2cm,top=1.6cm,bottom=1.8cm]{geometry}
\usepackage{fontspec}
\usepackage[russian]{babel}
\usepackage{amsmath}
\usepackage{unicode-math}
\setmainfont{Liberation Serif}
\setmonofont{Liberation Mono}[Scale=MatchLowercase]
\setmathfont{TeX Gyre Termes Math}
\usepackage{titlesec}
\usepackage{enumitem}
\usepackage{microtype}
\usepackage{booktabs}
\usepackage{caption}
\usepackage{float}
\usepackage{tikz}
\usepackage{subcaption}

\titleformat{\section}{\large\bfseries}{\thesection.}{0.5em}{}
\titleformat{\subsection}{\normalsize\bfseries}{\thesubsection.}{0.5em}{}
\titlespacing*{\section}{0pt}{1.0ex plus .3ex}{0.5ex}
\titlespacing*{\subsection}{0pt}{0.8ex plus .3ex}{0.4ex}
\setlength{\parindent}{0pt}
\setlength{\parskip}{0.45em}
\pagestyle{plain}
\captionsetup{font=small,labelsep=period,skip=4pt}

\definecolor{winc}{HTML}{2563EB}
\definecolor{visc}{HTML}{16A34A}
\definecolor{origc}{HTML}{94A3B8}
\definecolor{rejc}{HTML}{E5484D}
\definecolor{gridc}{HTML}{B8BDCC}

\usepackage[hidelinks,pdftitle={Лабораторная работа № 4. Алгоритмы отсечения отрезков},pdfauthor={Aliaksei Charnukha}]{hyperref}

\begin{document}

\begin{center}
  {\LARGE\bfseries Лабораторная работа № 4}\\[0.4em]
  {\large Алгоритмы отсечения отрезков}
\end{center}

\section{Цель}
Реализовать алгоритмы отсечения отрезков: Лианга~— Барски для прямоугольного окна
и Кируса~— Бека для выпуклого многоугольника. Визуализировать систему координат,
отсекатель, исходные и видимые части отрезков и сравнить время работы алгоритмов.

\section{Теория}

Отрезок задаётся параметрически:
\[
  P(t) = P_0 + t\,\mathbf{D}, \qquad \mathbf{D} = P_1 - P_0 = (\Delta x, \Delta y), \qquad t \in [0, 1],
\]
где $P_0 = (x_0, y_0)$, $P_1 = (x_1, y_1)$, $\Delta x = x_1 - x_0$, $\Delta y = y_1 - y_0$.
Выпуклый отсекатель~— пересечение полуплоскостей, поэтому видимая часть отрезка
задаётся интервалом $[t_0, t_1] \subseteq [0, 1]$.

\subsection{Алгоритм Лианга~— Барски (прямоугольное окно)}
Для окна $[x_{\min}, x_{\max}] \times [y_{\min}, y_{\max}]$ условие принадлежности
точки $P(t)$ окну записывается четырьмя линейными неравенствами $p_k\, t \le q_k$:
\[
\begin{aligned}
  k=1\text{ (левая)}:&\quad p_1 = -\Delta x,\; q_1 = x_0 - x_{\min}; &\qquad
  k=2\text{ (правая)}:&\quad p_2 = \Delta x,\; q_2 = x_{\max} - x_0;\\
  k=3\text{ (нижняя)}:&\quad p_3 = -\Delta y,\; q_3 = y_0 - y_{\min}; &\qquad
  k=4\text{ (верхняя)}:&\quad p_4 = \Delta y,\; q_4 = y_{\max} - y_0.
\end{aligned}
\]
Начальный интервал $[t_0, t_1] = [0, 1]$. Для каждого $k = 1, \dots, 4$:
\begin{itemize}[label=--, leftmargin=1.5em, itemsep=0.1em, topsep=0.2em]
  \item $p_k = 0$: отрезок параллелен границе; если $q_k < 0$, он вне окна и отвергается;
  \item $p_k < 0$ (вход в полуплоскость): $t_0 = \max(t_0,\, q_k / p_k)$;
  \item $p_k > 0$ (выход из полуплоскости): $t_1 = \min(t_1,\, q_k / p_k)$.
\end{itemize}
Если $t_0 > t_1$, отрезок невидим; иначе видимая часть~— $P(t_0)P(t_1)$.

Сквозной пример: окно $[1, 7] \times [2, 6]$, отрезок $P_0(-2, 1)$--$P_1(10, 7)$,
$\Delta x = 12$, $\Delta y = 6$.

\begin{table}[H]
\centering\small
\caption{Алгоритм Лианга~— Барски для отрезка $(-2,1)$--$(10,7)$ и окна $[1,7]\times[2,6]$.}
\begin{tabular}{clccclcc}
\toprule
$k$ & Граница & $p_k$ & $q_k$ & $r_k = q_k/p_k$ & Обновление & $t_0$ & $t_1$ \\
\midrule
1 & Левая ($x \ge 1$)   & $-12$ & $-3$ & $0.25$ & вход: $t_0 = \max(0, 0.25)$ & $0.25$ & $1$ \\
2 & Правая ($x \le 7$)  & $12$  & $9$  & $0.75$ & выход: $t_1 = \min(1, 0.75)$ & $0.25$ & $0.75$ \\
3 & Нижняя ($y \ge 2$)  & $-6$  & $-1$ & $0.167$ & вход: $t_0$ без изм. & $0.25$ & $0.75$ \\
4 & Верхняя ($y \le 6$) & $6$   & $5$  & $0.833$ & выход: $t_1$ без изм. & $0.25$ & $0.75$ \\
\bottomrule
\end{tabular}
\end{table}
Видимый отрезок: $P(0.25) = (1,\, 2.5)$, $P(0.75) = (7,\, 5.5)$ (рис.~\ref{fig:examples}, слева).

\subsection{Алгоритм Кируса~— Бека (выпуклый многоугольник)}
Пусть выпуклый многоугольник задан вершинами $V_0, \dots, V_{m-1}$, а $\mathbf{n}_i$~—
внутренняя нормаль к ребру $V_i V_{i+1}$. При обходе против часовой стрелки
для вектора ребра $\mathbf{e}_i = (e_x, e_y) = V_{i+1} - V_i$ внутренняя нормаль
равна $(-e_y, e_x)$; при обратном обходе знак меняется на $(e_y, -e_x)$.
Точка $P(t)$ лежит внутри полуплоскости ребра $i$, если
\[
  \mathbf{n}_i \cdot (P(t) - V_i) \ge 0
  \;\Longleftrightarrow\;
  d_i\, t + w_i \ge 0,
  \qquad d_i = \mathbf{n}_i \cdot \mathbf{D},\quad w_i = \mathbf{n}_i \cdot (P_0 - V_i).
\]
Начальный интервал $[t_0, t_1] = [0, 1]$. Для каждого ребра $i = 0, \dots, m - 1$:
\begin{itemize}[label=--, leftmargin=1.5em, itemsep=0.1em, topsep=0.2em]
  \item $d_i = 0$: отрезок параллелен ребру; при $w_i < 0$ он снаружи и отвергается;
  \item $d_i > 0$ (вход): $t_0 = \max(t_0,\, -w_i / d_i)$;
  \item $d_i < 0$ (выход): $t_1 = \min(t_1,\, -w_i / d_i)$.
\end{itemize}
При $t_0 > t_1$ отрезок невидим.

Сквозной пример: пятиугольник $(1,1), (6,1), (8,4), (5,7), (1,5)$,
отрезок $P_0(-1,1)$--$P_1(9,6)$, $\mathbf{D} = (10, 5)$.

\begin{table}[H]
\centering\small
\caption{Алгоритм Кируса~— Бека для отрезка $(-1,1)$--$(9,6)$ и выпуклого пятиугольника.}
\begin{tabular}{ccccccccc}
\toprule
Ребро & $\mathbf{n}_i$ & $w_i$ & $d_i$ & $-w_i/d_i$ & Тип & $t_0$ & $t_1$ \\
\midrule
1: $(1,1)$--$(6,1)$ & $(0, 5)$   & $0$   & $25$  & $0$     & вход  & $0$   & $1$ \\
2: $(6,1)$--$(8,4)$ & $(-3, 2)$  & $21$  & $-20$ & $1.05$  & выход & $0$   & $1$ \\
3: $(8,4)$--$(5,7)$ & $(-3, -3)$ & $36$  & $-45$ & $0.8$   & выход & $0$   & $0.8$ \\
4: $(5,7)$--$(1,5)$ & $(2, -4)$  & $12$  & $0$   & ---     & парал., $w_4 > 0$ & $0$ & $0.8$ \\
5: $(1,5)$--$(1,1)$ & $(4, 0)$   & $-8$  & $40$  & $0.2$   & вход  & $0.2$ & $0.8$ \\
\bottomrule
\end{tabular}
\end{table}
Видимый отрезок: $P(0.2) = (1, 2)$, $P(0.8) = (7, 5)$ (рис.~\ref{fig:examples}, справа).

\begin{figure}[H]
\centering
\begin{subfigure}[b]{0.48\linewidth}\centering\resizebox{\linewidth}{!}{\begin{tikzpicture}[scale=0.48]
  \draw[step=1,gridc!55,very thin] (-2.5,-0.5) grid (10.5,7.5);
  \draw[->,black!60] (-2.5,0) -- (10.7,0) node[right]{\scriptsize $x$};
  \draw[->,black!60] (0,-0.5) -- (0,7.7) node[above]{\scriptsize $y$};
  \fill[winc,opacity=0.10] (1,2) rectangle (7,6);
  \draw[winc,thick] (1,2) rectangle (7,6);
  \draw[origc,dashed,thick] (-2,1) -- (10,7);
  \draw[visc,ultra thick] (1,2.5) -- (7,5.5);
  \fill[visc] (1,2.5) circle (3.2pt);
  \fill[visc] (7,5.5) circle (3.2pt);
  \fill[black!70] (-2,1) circle (2.5pt) node[below left]{\scriptsize $P_0$};
  \fill[black!70] (10,7) circle (2.5pt) node[above right]{\scriptsize $P_1$};
  \draw[origc,dashed,thick] (2,3) -- (5,5);
  \draw[visc,ultra thick] (2,3) -- (5,5);
  \fill[visc] (2,3) circle (3.2pt);
  \fill[visc] (5,5) circle (3.2pt);
  \draw[origc,dashed,thick] (4,0) -- (9,4);
  \draw[visc,ultra thick] (6.5,2) -- (7,2.4);
  \fill[visc] (6.5,2) circle (3.2pt);
  \fill[visc] (7,2.4) circle (3.2pt);
  \draw[rejc,dashed,thick] (0,7) -- (8,7);
  \draw[rejc,dashed,thick] (7.5,1) -- (10,5);
\end{tikzpicture}
}\caption{Лианг~— Барски (окно $[1,7]\times[2,6]$)}\end{subfigure}\hfill
\begin{subfigure}[b]{0.48\linewidth}\centering\resizebox{\linewidth}{!}{\begin{tikzpicture}[scale=0.48]
  \draw[step=1,gridc!55,very thin] (-1.5,-0.5) grid (9.5,7.5);
  \draw[->,black!60] (-1.5,0) -- (9.7,0) node[right]{\scriptsize $x$};
  \draw[->,black!60] (0,-0.5) -- (0,7.7) node[above]{\scriptsize $y$};
  \fill[winc,opacity=0.10] (1,1) -- (6,1) -- (8,4) -- (5,7) -- (1,5) -- cycle;
  \draw[winc,thick] (1,1) -- (6,1) -- (8,4) -- (5,7) -- (1,5) -- cycle;
  \draw[->,winc!80,thin] (3.5,1) -- (3.5,1.75);
  \draw[->,winc!80,thin] (7,2.5) -- (6.376,2.916);
  \draw[->,winc!80,thin] (6.5,5.5) -- (5.97,4.97);
  \draw[->,winc!80,thin] (3,6) -- (3.335,5.329);
  \draw[->,winc!80,thin] (1,3) -- (1.75,3);
  \draw[origc,dashed,thick] (-1,1) -- (9,6);
  \draw[visc,ultra thick] (1,2) -- (7,5);
  \fill[visc] (1,2) circle (3.2pt);
  \fill[visc] (7,5) circle (3.2pt);
  \fill[black!70] (-1,1) circle (2.5pt) node[below left]{\scriptsize $P_0$};
  \fill[black!70] (9,6) circle (2.5pt) node[above right]{\scriptsize $P_1$};
  \draw[origc,dashed,thick] (2,2) -- (5,5);
  \draw[visc,ultra thick] (2,2) -- (5,5);
  \fill[visc] (2,2) circle (3.2pt);
  \fill[visc] (5,5) circle (3.2pt);
  \draw[origc,dashed,thick] (0,6.5) -- (8,6.5);
  \draw[visc,ultra thick] (4,6.5) -- (5.5,6.5);
  \fill[visc] (4,6.5) circle (3.2pt);
  \fill[visc] (5.5,6.5) circle (3.2pt);
  \draw[origc,dashed,thick] (3,0) -- (8,2);
  \draw[visc,ultra thick] (5.5,1) -- (6.182,1.273);
  \fill[visc] (5.5,1) circle (3.2pt);
  \fill[visc] (6.182,1.273) circle (3.2pt);
\end{tikzpicture}
}\caption{Кирус~— Бек (выпуклый пятиугольник и нормали)}\end{subfigure}
\caption{Синее~— отсекатель, пунктир~— исходный отрезок (красный при отказе), зелёное~— видимая часть.}
\label{fig:examples}
\end{figure}

\section{Ход работы}
Приложение написано на HTML, CSS и JavaScript и открывается в браузере без
установки (\texttt{lab4/index.html}). Алгоритмы и парсер входных данных вынесены
в модуль \texttt{clipping.js}. Реализовано:
\begin{itemize}[label=--, leftmargin=1.5em, itemsep=0.1em, topsep=0.2em]
  \item загрузка файла формата $n$, $n$ строк отрезков $x_1\ y_1\ x_2\ y_2$,
        координаты окна (или вершин многоугольника) и текстовое редактирование;
  \item масштабируемая и перетаскиваемая координатная сетка с осями, цветовым
        выделением отсекателя, исходных, отсечённых и отвергнутых отрезков;
  \item интерактивное добавление отрезков кликами, проверка выпуклости
        многоугольника и поддержка произвольного порядка обхода вершин (CW/CCW);
  \item таблица промежуточных вычислений для выбранного отрезка и замер времени.
\end{itemize}
Модуль проверен 8 юнит-тестами (\texttt{clipping.test.cjs}): сквозные примеры,
отрезки на границе, вырожденные отрезки, инвариантность к направлению обхода
многоугольника и совпадение Лианга~— Барски с Кирусом~— Беком на 500 случайных
отрезках.

\section{Эксперименты}
Набор~— $50\,000$ случайных отрезков с концами в квадрате $[-5, 15]^2$.
Для прямоугольного окна $[1, 7] \times [2, 6]$ полностью внутри оказалось
$0.4\,\%$ отрезков, отсечено $38.6\,\%$, отвергнуто $61.0\,\%$; расхождений между
Лиангом~— Барски и Кирусом~— Беком нет.

Замер времени~— лучший из 12 прогонов после прогрева JIT (Node.js~20):

\begin{table}[H]
\centering\small
\caption{Время отсечения на один отрезок, нс.}
\begin{tabular}{lc}
\toprule
Алгоритм / число сторон $m$ & Время, нс/отрезок \\
\midrule
Лианг~— Барски (прямоугольник) & 652 \\
Кирус~— Бек ($m = 4$, прямоугольник) & 739 \\
Кирус~— Бек ($m = 8$) & 1151 \\
Кирус~— Бек ($m = 16$) & 2160 \\
Кирус~— Бек ($m = 32$) & 4043 \\
\bottomrule
\end{tabular}
\end{table}

\section{Результаты}
\begin{itemize}[label=--, leftmargin=1.5em, itemsep=0.2em, topsep=0.2em]
  \item Разработано веб-приложение с алгоритмами Лианга~— Барски и Кируса~— Бека,
        поддержкой входного файла и пошаговой таблицей; пройдено 8 юнит-тестов.
  \item На $50\,000$ случайных отрезках оба алгоритма на прямоугольном окне дали
        полностью совпадающие результаты.
  \item Среднее время на отрезок: Лианг~— Барски~— 652~нс, Кирус~— Бек для $m = 4$~—
        739~нс, при росте числа вершин от 4 до 32 время растёт с 739 до 4043~нс.
\end{itemize}

\section{Выводы}
\begin{itemize}[label=--, leftmargin=1.5em, itemsep=0.2em, topsep=0.2em]
  \item Оба алгоритма используют параметрическое сужение интервала $[t_0, t_1]$
        полуплоскостями: Лианг~— Барски является специализацией схемы Кируса~— Бека
        для границ, параллельных осям координат.
  \item Специализация в Лианге~— Барски избавляет от вычисления нормалей и скалярных
        произведений, поэтому на прямоугольном окне он примерно на $12\,\%$ быстрее
        Кируса~— Бека при идентичном результате.
  \item Трудоёмкость Кируса~— Бека линейно зависит от числа вершин отсекателя $O(m)$,
        что подтверждается ростом времени почти вдвое при каждом удвоении $m$.
\end{itemize}

\end{document}
